\documentclass{article}
\usepackage{amsmath,amssymb,amsthm}
\usepackage{algorithm}
\usepackage{algpseudocode}
\usepackage{hyperref}
\newtheorem{theorem}{Theorem}
\newtheorem{lemma}{Lemma}
\newtheorem{assumption}{Assumption}
\newtheorem{corollary}{Corollary}
\newcommand{\prox}{\operatorname{prox}}
\newcommand{\R}{\mathbb{R}}

\begin{document}

\section*{Proximal Gradient Method}

\section*{Mathematical Formulation}
Let 
\[
F(x)\;:=\;f(x)+g(x),\qquad x\in\R^{d},
\]
where  

\begin{itemize}
    \item $f:\R^{d}\to\R$ is convex, differentiable and its gradient is $L$–Lipschitz, i.e.
    \[
        \|\nabla f(u)-\nabla f(v)\|\le L\|u-v\|,\qquad\forall u,v\in\R^{d};
    \]
    \item $g:\R^{d}\to\R\cup\{+\infty\}$ is proper, closed and convex (possibly non‑smooth);
    \item $\prox_{\eta g}:\R^{d}\to\R^{d}$ denotes the proximal operator
    \[
        \prox_{\eta g}(y)\;:=\;\arg\min_{z\in\R^{d}}\Big\{ g(z)+\frac{1}{2\eta}\|z-y\|^{2}\Big\},
    \]
    which is well defined for every $\eta>0$.
\end{itemize}

Given a stepsize $\eta>0$, the \emph{proximal gradient iteration} reads
\begin{equation}
\label{eq:pg}
x_{k+1}\;=\;\prox_{\eta g}\bigl(x_{k}-\eta\nabla f(x_{k})\bigr),\qquad k=0,1,2,\dots
\end{equation}
When $g\equiv0$, $\prox_{\eta g}$ reduces to the identity map and \eqref{eq:pg} becomes ordinary gradient descent.

\section*{Pseudocode}
\begin{algorithm}[H]
\caption{Proximal Gradient Method}
\begin{algorithmic}[1]
\Require Convex $f$, convex $g$, stepsize $\eta\in(0,1/L]$, initial point $x_{0}\in\R^{d}$, max iterations $K$ (or tolerance $\varepsilon$).
\For{$k=0$ \textbf{to} $K-1$}
    \State Compute gradient $g_{k}\gets \nabla f(x_{k})$.
    \State Take a forward step $y_{k}\gets x_{k}-\eta g_{k}$.
    \State Apply the proximal map $x_{k+1}\gets \prox_{\eta g}(y_{k})$.
    \If{$\|x_{k+1}-x_{k}\|\le\varepsilon$}
        \State \textbf{break}
    \EndIf
\EndFor
\State \Return $x_{K}$ (or the last iterate).
\end{algorithmic}
\end{algorithm}

\section*{Dry Run Example}
Consider the one‑dimensional problem
\[
f(w)=(w-3)^{2},\qquad g\equiv0,
\]
so that $\prox_{\eta g}(y)=y$.  
Take $w_{0}=0$, stepsize $\eta=0.1$, and run three iterations.

\begin{center}
\begin{tabular}{c|c|c|c}
Iteration $k$ & $w_{k}$ & $\nabla f(w_{k})=2(w_{k}-3)$ & $w_{k+1}=w_{k}-\eta\nabla f(w_{k})$\\\hline
0 & $0$   & $-6$ & $0-0.1(-6)=0.6$\\
1 & $0.6$ & $-4.8$ & $0.6-0.1(-4.8)=1.08$\\
2 & $1.08$& $-3.84$& $1.08-0.1(-3.84)=1.464$\\
\end{tabular}
\end{center}

Objective values:
\[
F(w_{0})=9,\quad F(w_{1})=(0.6-3)^{2}=5.76,\quad 
F(w_{2})=(1.08-3)^{2}=3.6864,\quad
F(w_{3})=(1.464-3)^{2}=2.3617.
\]

The iterates move monotonically toward the minimiser $w^{*}=3$.

\newpage

\section*{Assumptions}
\begin{assumption}[Problem regularity]\label{as:regular}
\begin{enumerate}
    \item $f$ is convex and differentiable with $L$‑Lipschitz continuous gradient.
    \item $g$ is proper, closed and convex.
    \item The composite function $F=f+g$ admits at least one minimiser $x^{*}\in\arg\min_{x}F(x)$.
    \item The stepsize satisfies $0<\eta\le 1/L$.
\end{enumerate}
\end{assumption}

\section*{Main Convergence Theorem}
\begin{theorem}[Ergodic $O(1/k)$ rate]\label{thm:rate}
Under Assumption~\ref{as:regular}, the iterates generated by \eqref{eq:pg} satisfy for every $k\ge 0$
\[
F(x_{k})-F(x^{*})\;\le\;\frac{1}{2\eta\,k}\,\|x_{0}-x^{*}\|^{2}.
\]
Consequently,
\[
\min_{0\le i\le k-1}\bigl(F(x_{i})-F(x^{*})\bigr)\;\le\;\frac{1}{2\eta\,k}\,\|x_{0}-x^{*}\|^{2},
\]
i.e. the method attains a sublinear convergence rate $O(1/k)$ in function value.
\end{theorem}

\section*{Theoretical Properties}
\begin{itemize}
    \item \textbf{Stability.} The algorithm is \emph{non‑expansive} in the sense that the proximal step does not increase distances:
    \[
        \|\prox_{\eta g}(u)-\prox_{\eta g}(v)\|\le\|u-v\|,\qquad\forall u,v\in\R^{d},
    \]
    which is the key ingredient of the proof.
    \item \textbf{Hyperparameter sensitivity.} The stepsize must respect $\eta\le 1/L$. Larger $\eta$ may violate the descent property, while smaller $\eta$ slows down convergence (the bound deteriorates linearly with $1/\eta$).
    \item \textbf{Comparison to baselines.}
    \begin{enumerate}
        \item With $g\equiv0$, Theorem~\ref{thm:rate} reduces to the classical $O(1/k)$ rate of gradient descent for convex smooth functions.
        \item Compared with subgradient methods for non‑smooth $g$, the proximal gradient method exploits the structure of $g$ and often yields dramatically faster practical performance while preserving the same theoretical rate.
    \end{enumerate}
\end{itemize}

\section*{Discussion}
The $O(1/k)$ bound is optimal for first‑order methods on the class of convex, Lipschitz‑smooth problems \cite{nesterov1983}. Accelerated variants (e.g., FISTA) can improve the rate to $O(1/k^{2})$ by incorporating momentum, but the basic proximal gradient method already provides a clean and robust baseline for composite optimisation.

\newpage

\section*{Preliminaries (Definitions and Lemmas)}
\begin{lemma}[Descent Lemma]\label{lem:descent}
If $\nabla f$ is $L$‑Lipschitz, then for any $u,v\in\R^{d}$
\[
f(v)\le f(u)+\langle\nabla f(u),v-u\rangle+\frac{L}{2}\|v-u\|^{2}.
\]
\end{lemma}
\begin{proof}
Standard; follows from integrating the Lipschitz condition. \hfill $\square$
\end{proof}

\begin{lemma}[Optimality condition of the proximal map]\label{lem:prox-opt}
For any $y\in\R^{d}$ and $\eta>0$, $z^{+}=\prox_{\eta g}(y)$ satisfies
\[
0\in \partial g(z^{+})+\frac{1}{\eta}(z^{+}-y).
\]
Equivalently,
\[
\langle z-z^{+},\,\frac{1}{\eta}(y-z^{+})\rangle + g(z)-g(z^{+})\ge 0,\qquad\forall z\in\R^{d}.
\]
\end{lemma}
\begin{proof}
By definition $z^{+}$ minimises $g(z)+\frac{1}{2\eta}\|z-y\|^{2}$, thus the first‑order optimality condition yields the inclusion. \hfill $\square$
\end{proof}

\begin{lemma}[Non‑expansiveness of the proximal operator]\label{lem:nonexp}
For any $u,v\in\R^{d}$,
\[
\|\prox_{\eta g}(u)-\prox_{\eta g}(v)\|\le\|u-v\|.
\]
\end{lemma}
\begin{proof}
Direct consequence of the firm non‑expansiveness property of proximal maps. \hfill $\square$
\end{proof}

\newpage

\section*{Main Proof (Step‑by‑Step Derivation)}
We prove Theorem~\ref{thm:rate}. Throughout the proof, fix an arbitrary minimiser $x^{*}\in\arg\min F$.

\subsection*{Step 1 – Apply optimality of the proximal step}
From Lemma~\ref{lem:prox-opt} with $y_{k}=x_{k}-\eta\nabla f(x_{k})$ and $z^{+}=x_{k+1}$ we obtain for any $z\in\R^{d}$:
\[
\langle z-x_{k+1},\; \nabla f(x_{k}) +\tfrac{1}{\eta}(x_{k+1}-x_{k})\rangle + g(z)-g(x_{k+1})\ge 0.
\tag{1}
\label{eq:opt}
\]
Choose $z=x^{*}$. Rearranging \eqref{eq:opt} gives
\[
\langle x^{*}-x_{k+1},\nabla f(x_{k})\rangle
\le \frac{1}{\eta}\langle x^{*}-x_{k+1},x_{k}-x_{k+1}\rangle + g(x^{*})-g(x_{k+1}).
\tag{2}
\label{eq:opt-star}
\]

\subsection*{Step 2 – Expand the inner product}
Using the identity $2\langle a-b,c-b\rangle = \|a-b\|^{2}+\|c-b\|^{2}-\|a-c\|^{2}$ with $a=x^{*}$, $c=x_{k}$, $b=x_{k+1}$,
\[
2\langle x^{*}-x_{k+1},x_{k}-x_{k+1}\rangle
= \|x^{*}-x_{k+1}\|^{2}+\|x_{k}-x_{k+1}\|^{2}-\|x^{*}-x_{k}\|^{2}.
\tag{3}
\label{eq:inner}
\]
Dividing \eqref{eq:inner} by $2\eta$ and inserting into \eqref{eq:opt-star} yields
\[
\langle x^{*}-x_{k+1},\nabla f(x_{k})\rangle
\le \frac{1}{2\eta}\bigl(\|x^{*}-x_{k+1}\|^{2}+\|x_{k}-x_{k+1}\|^{2}-\|x^{*}-x_{k}\|^{2}\bigr)
+ g(x^{*})-g(x_{k+1}).
\tag{4}
\label{eq:opt-expanded}
\]

\subsection*{Step 3 – Use convexity of $f$}
Convexity of $f$ gives
\[
f(x_{k+1})\ge f(x_{k})+\langle\nabla f(x_{k}),x_{k+1}-x_{k}\rangle.
\tag{5}
\label{eq:convex-f}
\]
Equivalently,
\[
\langle x^{*}-x_{k+1},\nabla f(x_{k})\rangle
= \langle x^{*}-x_{k},\nabla f(x_{k})\rangle + \langle x_{k}-x_{k+1},\nabla f(x_{k})\rangle.
\tag{6}
\label{eq:split}
\]

\subsection*{Step 4 – Combine \eqref{eq:opt-expanded} and \eqref{eq:split}}
Insert \eqref{eq:split} into the left‑hand side of \eqref{eq:opt-expanded} and rearrange:
\[
\langle x^{*}-x_{k},\nabla f(x_{k})\rangle
\le \frac{1}{2\eta}\bigl(\|x^{*}-x_{k+1}\|^{2}+\|x_{k}-x_{k+1}\|^{2}-\|x^{*}-x_{k}\|^{2}\bigr)
+ g(x^{*})-g(x_{k+1}) - \langle x_{k}-x_{k+1},\nabla f(x_{k})\rangle .
\tag{7}
\label{eq:combined}
\]

\subsection*{Step 5 – Apply the descent lemma to control $\langle x_{k}-x_{k+1},\nabla f(x_{k})\rangle$}
From Lemma~\ref{lem:descent} with $u=x_{k}$, $v=x_{k+1}$,
\[
f(x_{k+1})\le f(x_{k})+\langle\nabla f(x_{k}),x_{k+1}-x_{k}\rangle+\frac{L}{2}\|x_{k+1}-x_{k}\|^{2}.
\]
Rearranging,
\[
-\langle x_{k}-x_{k+1},\nabla f(x_{k})\rangle
\le f(x_{k})-f(x_{k+1})+\frac{L}{2}\|x_{k+1}-x_{k}\|^{2}.
\tag{8}
\label{eq:descent}
\]

\subsection*{Step 6 – Substitute \eqref{eq:descent} into \eqref{eq:combined}}
\[
\langle x^{*}-x_{k},\nabla f(x_{k})\rangle
\le \frac{1}{2\eta}\bigl(\|x^{*}-x_{k+1}\|^{2}+\|x_{k}-x_{k+1}\|^{2}-\|x^{*}-x_{k}\|^{2}\bigr)
+ g(x^{*})-g(x_{k+1})
\]
\[
\quad\; +\; f(x_{k})-f(x_{k+1})+\frac{L}{2}\|x_{k+1}-x_{k}\|^{2}.
\tag{9}
\label{eq:pre-final}
\]

\subsection*{Step 7 – Use convexity of $f$ again to bound $\langle x^{*}-x_{k},\nabla f(x_{k})\rangle$}
Convexity yields
\[
f(x_{k})-f(x^{*})\le \langle\nabla f(x_{k}),x_{k}-x^{*}\rangle
= -\langle x^{*}-x_{k},\nabla f(x_{k})\rangle .
\tag{10}
\label{eq:convex-f-star}
\]
Thus,
\[
-\langle x^{*}-x_{k},\nabla f(x_{k})\rangle \ge f(x_{k})-f(x^{*}).
\tag{11}
\label{eq:lower}
\]

\subsection*{Step 8 – Plug \eqref{eq:lower} into \eqref{eq:pre-final}}
Multiplying \eqref{eq:pre-final} by $-1$ and using \eqref{eq:lower} gives
\[
f(x_{k})-f(x^{*}) + g(x_{k+1})-g(x^{*})
\le \frac{1}{2\eta}\bigl(\|x^{*}-x_{k}\|^{2}-\|x^{*}-x_{k+1}\|^{2}\bigr)
+\Bigl(\frac{L}{2}-\frac{1}{2\eta}\Bigr)\|x_{k+1}-x_{k}\|^{2}.
\tag{12}
\label{eq:intermediate}
\]

\subsection*{Step 9 – Choose stepsize $\eta\le 1/L$}
Since $0<\eta\le 1/L$, we have $L/2-1/(2\eta)\le 0$, therefore the last term in \eqref{eq:intermediate} is non‑positive and can be dropped:
\[
F(x_{k+1})-F(x^{*})
\le \frac{1}{2\eta}\bigl(\|x^{*}-x_{k}\|^{2}-\|x^{*}-x_{k+1}\|^{2}\bigr).
\tag{13}
\label{eq:descent-final}
\]
\textbf{[Step 13]} This is the key descent inequality.

\subsection*{Step 10 – Telescoping sum}
Sum \eqref{eq:descent-final} from $k=0$ to $k=K-1$:
\[
\sum_{k=0}^{K-1}\bigl(F(x_{k+1})-F(x^{*})\bigr)
\le \frac{1}{2\eta}\bigl(\|x^{*}-x_{0}\|^{2}-\|x^{*}-x_{K}\|^{2}\bigr)
\le \frac{1}{2\eta}\|x_{0}-x^{*}\|^{2}.
\tag{14}
\label{eq:sum}
\]

\subsection*{Step 11 – Derive the $O(1/k)$ bound}
Since the left‑hand side of \eqref{eq:sum} contains $K$ non‑negative terms,
\[
K\;\min_{0\le i\le K-1}\bigl(F(x_{i+1})-F(x^{*})\bigr)
\le \frac{1}{2\eta}\|x_{0}-x^{*}\|^{2}.
\]
Consequently,
\[
\min_{0\le i\le K-1}\bigl(F(x_{i+1})-F(x^{*})\bigr)
\le \frac{1}{2\eta K}\|x_{0}-x^{*}\|^{2}.
\tag{15}
\label{eq:rate}
\]
Because $F(x_{k})$ is non‑increasing by \eqref{eq:descent-final}, the same bound holds for each individual iterate:
\[
F(x_{k})-F(x^{*})\le\frac{1}{2\eta k}\|x_{0}-x^{*}\|^{2},\qquad k\ge1.
\]
This completes the proof of Theorem~\ref{thm:rate}. \hfill $\square$

\section*{Rate Derivation (With Constants)}
The constant in the bound is explicitly $C=\frac{1}{2\eta}\|x_{0}-x^{*}\|^{2}$. If the stepsize is chosen as the largest admissible value $\eta=1/L$, the rate becomes
\[
F(x_{k})-F(x^{*})\le \frac{L}{2k}\,\|x_{0}-x^{*}\|^{2}.
\]

\section*{Special Cases}
\begin{itemize}
    \item \textbf{Pure gradient descent ($g\equiv0$).} The proximal map reduces to identity, and the proof simplifies to the classic $O(1/k)$ rate for smooth convex functions.
    \item \textbf{Indicator function $g=\iota_{C}$ of a closed convex set $C$.} Then $\prox_{\eta g}$ is the Euclidean projection onto $C$, and the method becomes projected gradient descent with the same $O(1/k)$ guarantee.
    \item \textbf{Lasso regularisation $g(x)=\lambda\|x\|_{1}$.} The prox is the soft‑thresholding operator; the bound remains unchanged, showing that sparsity‑inducing regularisers do not deteriorate the convergence order.
\end{itemize}

\begin{thebibliography}{9}
\bibitem{nesterov1983}
Y.~Nesterov.
\newblock {\em A method of solving a convex programming problem with convergence rate $O(1/k^{2})$}.
\newblock Soviet Mathematics Doklady, 27(2):372–376, 1983.
\end{thebibliography}

\end{document}